\documentclass[10pt,a4paper]{amsart}
\usepackage{amsmath,amssymb,amsthm,graphicx}
\usepackage{amscd}
\usepackage{hyperref}
\numberwithin{equation}{section}

\newtheorem{Theorem}{Theorem}[section]
\newtheorem{Corollary}[Theorem]{Corollary}
\newtheorem{Lemma}[Theorem]{Lemma}
\newtheorem{Proposition}[Theorem]{Proposition}
{\theoremstyle{definition}
\newtheorem{Note}[Theorem]{Note}
\newtheorem{Definition}[Theorem]{Definition}
\newtheorem{Problem}[Theorem]{Problem}
}
\usepackage[a4paper,margin=22mm]{geometry}
\title{Faithfulness of the modular-group action on the universal Askey–Wilson algebra}
\author{Paul Terwilliger}
\date{doi:10.3842/SIGMA.2011.069; CC BY 4.0}
\begin{document}
\maketitle
\noindent\textit{Allowed standard prerequisites (editorial).}\begin{itemize}
\item PSL2(Z)=C3*C2 with reduced alternating-word normal form. Alperin1993, original bibliography alpern. Representation construction and nontrivial reduced-word test.
\item Center of Mat2(F[lambda,lambda\textasciicircum{}-1]) consists of scalar matrices. Standard ring-theoretic fact; author statement and short proof marker retained at628–642.
\item Basic polynomial degree, determinant and matrix multiplication. Leading coefficient and determinant contradiction.
\end{itemize}\medskip
Our conventions for the paper are as follows.
An algebra is meant to be associative and have a 1.
A subalgebra has the same 1 as the parent algebra.
We f\/ix a f\/ield $\mathbb F$
and  a nonzero $q \in \mathbb F$ such that $q^4\not=1$.

\begin{Definition}
\label{def:uaw}
Def\/ine an $\mathbb F$-algebra
$\Delta=\Delta_q$
 by generators and relations in the following way.
The generators are~$A$,~$B$,~$C$.
The relations assert that each of
\begin{gather}
 A+ \frac{qBC-q^{-1}CB}{q^2-q^{-2}},
%\label{eq:u1}
\qquad B+
\frac{qCA-q^{-1}AC}{q^2-q^{-2}},
%\label{eq:u2}
\qquad
C+
\frac{qAB-q^{-1}BA}{q^2-q^{-2}}
%\label{eq:u3}
\label{eq:comlist}
\end{gather}
is central in $\Delta$.
We call $\Delta$ the {\it universal Askey--Wilson algebra}.
\end{Definition}





\begin{Definition}
\label{def:abc}
\rm
For the
three central elements in
(\ref{eq:comlist}), multiply each by $q+q^{-1}$ to get~$\alpha$,~$\beta$,~$\gamma$. Thus
\begin{gather}
A+ \frac{qBC-q^{-1}CB}{q^2-q^{-2}}
=
\frac{\alpha}{q+q^{-1}},
\label{eq:u1}
\\
B+
\frac{qCA-q^{-1}AC}{q^2-q^{-2}}
=
\frac{\beta}{q+q^{-1}},
\label{eq:u2}
\\
C+
\frac{qAB-q^{-1}BA}{q^2-q^{-2}}
=
\frac{\gamma}{q+q^{-1}}.
\label{eq:u3}
\end{gather}
Note that each of $\alpha$, $\beta$, $\gamma$ is central in
$\Delta$.
(The purpose of the factor $q+q^{-1}$ is
to make the upcoming formula~(\ref{eq:caspreview})
more attractive.)
\end{Definition}

\begin{gather}
\label{eq:caspreview}
\Omega =
qABC+q^2A^2+q^{-2}B^2+q^2C^2-qA\alpha-q^{-1}B\beta -q C\gamma.
\end{gather}

\begin{Lemma}
\label{lem:abgpres}
The $\mathbb F$-algebra $\Delta $ is generated by
$A$, $B$, $\gamma$. Moreover
\begin{gather}
C  =  \frac{\gamma}{q+q^{-1}}
-\frac{qAB-q^{-1}BA}{q^2-q^{-2}},
\label{eq:getc}
\\
\alpha  =  \frac{B^2A-(q^2+q^{-2})BAB+AB^2+
(q^2-q^{-2})^2A+(q-q^{-1})^2B\gamma}{(q-q^{-1})(q^2-q^{-2})},
\label{eq:getalpha}
\\
\beta  =  \frac{A^2B-(q^2+q^{-2})ABA+BA^2+
(q^2-q^{-2})^2B+(q-q^{-1})^2A\gamma}{(q-q^{-1})(q^2-q^{-2})}.
\label{eq:getbeta}
\end{gather}
\end{Lemma}

\begin{proof}\sloppy
Line~(\ref{eq:getc})
is from (\ref{eq:u3}).
To get
(\ref{eq:getalpha}),
(\ref{eq:getbeta})
eliminate $C$ in
(\ref{eq:u1}),
(\ref{eq:u2})
using
(\ref{eq:getc}).
\end{proof}

 Recall the notation
\begin{gather*}
\lbrack n \rbrack_q= \frac{q^n-q^{-n}}{q-q^{-1}}, \qquad
n=0,1,2,\ldots.
\end{gather*}




{\samepage

\begin{Theorem}
\label{prop:altpres}
The $\mathbb F$-algebra $\Delta $ has a  presentation by generators
$A$, $B$, $\gamma$ and
relations
 \begin{gather*}
 A^3B-\lbrack 3\rbrack_q A^2BA+\lbrack 3\rbrack_q ABA^2-BA^3
= -\big(q^2-q^{-2}\big)^2(AB-BA),
\\
 B^3A-\lbrack 3\rbrack_q B^2AB+\lbrack 3\rbrack_q BAB^2-AB^3
= -\big(q^2-q^{-2}\big)^2(BA-AB),
\\
 A^2B^2-B^2A^2+\big(q^2+q^{-2}\big)(BABA-ABAB) = -\big(q-q^{-1}\big)^2(AB-BA)\gamma,
\\
  \qquad \gamma A = A \gamma,   \qquad \gamma B=B \gamma.
 \end{gather*}
\end{Theorem}

\begin{proof}
Use Lemma~\ref{lem:abgpres} to
express the def\/ining relations
for $\Delta $ in terms of $A$, $B$, $\gamma$.
\end{proof}}

We now consider
some automorphisms of~$\Delta $.
Recall that the modular group
${\rm  {PSL}}_2(\mathbb Z)$ has a~presentation
by generators~$\rho$, $\sigma$ and
relations $\rho^3=1$, $\sigma^2=1$.
See for example \cite{alpern}.


\begin{Theorem}
\label{thm:try}
The group
${\rm  {PSL}}_2(\mathbb Z)$ acts on
$\Delta $ as a group of automorphisms in the following way:
\[
\begin{array}{@{}c| ccc | c c c}
u &  A & B  & C
&  \alpha & \beta  & \gamma
\\
\hline
\rho(u) &  B & C & A
&  \beta  & \gamma & \alpha
\\
\sigma(u) &  B & A & C+\frac{AB-BA}{q-q^{-1}}
&  \beta & \alpha & \gamma
\end{array}
\]
\end{Theorem}

\begin{proof}
By Def\/inition
\ref{def:uaw} there exists an automorphism
$P$ of $\Delta $ that sends
\begin{gather*}
A \mapsto B,
\qquad
B \mapsto C,
\qquad
C \mapsto A.
\end{gather*}
Observe $P^3=1$. By
(\ref{eq:u1})--(\ref{eq:u3})
the map $P$ sends
\begin{gather*}
\alpha \mapsto \beta,
\qquad
\beta \mapsto \gamma,
\qquad
\gamma \mapsto \alpha.
\end{gather*}
By Theorem
\ref{prop:altpres} there exists
an automorphism
$S$ of $\Delta $ that sends
\begin{gather*}
A \mapsto B,
 \qquad
B \mapsto A,
\qquad
\gamma \mapsto \gamma.
\end{gather*}
Observe $S^2=1$. By
Lemma
\ref{lem:abgpres}
the map $S$
sends
\begin{gather*}
\alpha \mapsto \beta,
\qquad
\beta \mapsto \alpha,
\qquad
C \mapsto C + \frac{AB-BA}{q-q^{-1}}.
\end{gather*}
The result follows.
\end{proof}


In Theorem
\ref{thm:try}
we gave an
 action of ${\rm  {PSL}}_2(\mathbb Z)$  on $\Delta $.
Our next goal is to show that this
action is faithful.

Let $\lambda $ denote an indeterminate.
Let $\mathbb F\lbrack \lambda,\lambda^{-1}\rbrack$
denote the $\mathbb F$-algebra consisting of the Laurent
polynomials in $\lambda$ that have all coef\/f\/icients
in~$\mathbb F$.
We will be discussing the $\mathbb F$-algebra
\begin{gather*}
\Lambda ={\rm {Mat}}_2(\mathbb F) \otimes_{\mathbb F} \mathbb F\big\lbrack \lambda, \lambda^{-1}\big\rbrack.
\end{gather*}
We view elements of $\Lambda $ as $2\times 2$ matrices that have entries
in
$\mathbb F\lbrack \lambda, \lambda^{-1}\rbrack $.
From this point of view the product operation
for $\Lambda $ is ordinary matrix multiplication,
and the multiplicative identity in~$\Lambda $
is the identity matrix~$I$.
For notational convenience def\/ine
$\mu = \lambda + \lambda^{-1}$.


For later use we now describe the center $Z(\Lambda)$.{\samepage
\begin{Lemma}
\label{lem:zh}
For all $\eta \in \Lambda $ the following
$(i)$, $(ii)$
are equivalent:
\begin{enumerate}\itemsep=0pt
\item[$(i)$] $\eta  \in Z(\Lambda)$.
\item[$(ii)$] There exists $\theta \in
 \mathbb F\lbrack \lambda, \lambda^{-1} \rbrack $ such that $\eta=\theta I$.
\end{enumerate}
\end{Lemma}

\begin{proof}
Routine.
\end{proof}}



\begin{Definition}
\label{def:xyz}
Let $ \mathcal A$, $\mathcal B$, $\mathcal C$ denote the following elements of
$\Lambda $:
\begin{gather*}
 \mathcal A =
\left(
\begin{array}{  c c}
\lambda & 1-\lambda^{-1}   \\
0 & \lambda^{-1}
\end{array}
\right),
\qquad
\mathcal B =
\left(
\begin{array}{  c c}
\lambda^{-1} &  0 \\
\lambda-1 & \lambda
\end{array}
\right),
 \qquad \mathcal C =
\left(
\begin{array}{  c c}
1 &  1-\lambda \\
\lambda^{-1}-1 & \lambda+ \lambda^{-1}-1
\end{array}
\right).
\end{gather*}
\end{Definition}


\begin{Lemma}
\label{lem:xyz}
We have
\begin{gather*}
\mathcal A \mathcal B\mathcal C=I, \qquad
\mathcal A+{\mathcal A}^{-1}=\mu I,
\qquad
\mathcal B+{\mathcal B}^{-1}=\mu I,
\qquad
\mathcal C+{\mathcal C}^{-1}=\mu I.
\end{gather*}
\end{Lemma}

\begin{proof}
Use Def\/inition
\ref{def:xyz}.
\end{proof}

\begin{Lemma}
\label{lem:multtable2}
In the algebra $\Lambda $ the elements $\mathcal A$, $\mathcal B$, $\mathcal C$
multiply as follows:
\[
\begin{array}{@{}c| c c c}
 &  \mathcal A & \mathcal B  & \mathcal C \\
\hline
\mathcal A &   \mu \mathcal A  - I &   \mu I - \mathcal C
&  \mu \mathcal A  + \mathcal B + \mu \mathcal C  -\mu^2 I \\
\mathcal B
&
\mu \mathcal B + \mathcal C + \mu \mathcal A  -\mu^2 I
&
  \mu \mathcal B - I
&
 \mu I- \mathcal A
\\
\mathcal C
&
 \mu I - \mathcal B
&
\mu \mathcal C  + \mathcal A + \mu \mathcal B  -\mu^2 I
&
\mu \mathcal C - I
\end{array}
\]
\end{Lemma}

\begin{proof}
Use Lemma
\ref{lem:xyz}.
\end{proof}




 The algebra $\Lambda $ is not generated by
$\mathcal A$, $\mathcal B$, $\mathcal C$. However
we do have the following result.

\begin{Lemma}
\label{lem:whatittakes}
Suppose  $\eta  \in \Lambda $ commutes with at
 least two of
$\mathcal A$, $\mathcal B$, $\mathcal C$.
Then $\eta  \in Z(\Lambda )$.
\end{Lemma}

\begin{proof}
For $1 \leq i,j\leq 2$
let $\eta_{ij}$ denote the $(i,j)$-entry of $\eta$.
The matrix $\eta $ commutes with each of~$\mathcal A$,~$\mathcal B$,~$\mathcal C$
since $\mathcal A\mathcal B\mathcal C=I$.
In the equation $\eta \mathcal A=\mathcal A \eta  $,
evaluate $\mathcal A$ using Def\/inition~\ref{def:xyz}, and simplify the result to get
$\eta_{21}=0$. Similarly using $\eta \mathcal B=\mathcal B\eta $
we f\/ind $\eta_{12}=0$ and
$\eta_{11}=\eta_{22}$.
Therefore $\eta =\eta_{11} I\in Z(\Lambda )$.
\end{proof}


  Next we describe an
action of ${\rm  {PSL}}_2(\mathbb Z)$ on $\Lambda $ as
a group of automorphisms.

\begin{Definition}
\label{def:ps}
Let $p$ and $s$ denote the following elements
of $\Lambda $:
\begin{gather*}
p =
\left(
\begin{array}{  c c}
0 & -1   \\
1 & 1
\end{array}
\right),
\qquad
s =
\left(
\begin{array}{  c c}
0 &  1 \\
\lambda  & 0
\end{array}
\right).
\end{gather*}
\end{Definition}

\begin{Lemma}
\label{lem:4}
The following $(i)$--$(iv)$ hold.
\begin{enumerate}\itemsep=0pt
\item[$(i)$] $\det (p)=1$ and  $\det (s)=-\lambda$.
\item[$(ii)$]
 $p^3=-I$ and $s^2=\lambda I $.
\item[$(iii)$]
$p \mathcal A p^{-1} = \mathcal B$,
$p  \mathcal B p^{-1} = \mathcal C$,
$p \mathcal C p^{-1} = \mathcal A$.
\item[$(iv)$]
$s \mathcal A s^{-1} = \mathcal B$ and
$s \mathcal B s^{-1} = \mathcal A$.
\end{enumerate}
\end{Lemma}

\begin{proof}
$(i)$, $(ii)$
Use Def\/inition~\ref{def:ps}.
$(iii)$, $(iv)$ Use
Def\/inition~\ref{def:xyz}
and
Def\/inition~\ref{def:ps}.
\end{proof}





\begin{Lemma}
\label{lem:conj}
The group
${\rm  {PSL}}_2(\mathbb Z)$ acts on $\Lambda $ as a group of
automorphisms such
that
$\rho(\eta )= p \eta p^{-1}$
and
$\sigma (\eta ) = s \eta s^{-1}$ for all
$\eta  \in \Lambda $.
\end{Lemma}

\begin{proof}
By Lemma~\ref{lem:4}$(ii)$ the elements  $p^3$, $s^2$ are in $Z(\Lambda)$.
\end{proof}




\begin{Lemma}
\label{lem:faith1}
The action of
${\rm  {PSL}}_2(\mathbb Z)$ on $\Lambda $ is faithful.
\end{Lemma}

\begin{proof}
Pick an integer $n\geq 1$.
Consider an element $\eta \in \Lambda $
of the form
$\eta =\eta_1 \eta_2 \cdots \eta_n$ such that
 for $1 \leq i \leq n$,  $\eta_i=s$
for one parity of $i$ and
$\eta_i \in \lbrace p,p^{-1} \rbrace$ for the other parity of~$i$.
We show that $\eta \not\in Z(\Lambda)$.
To this end we assume
 $\eta \in Z(\Lambda )$ and get a contradiction.
Let $\ell$ denote the number of times~$s$ occurs among
$\lbrace \eta_i\rbrace_{i=1}^n$.
Assume for the moment $\ell = 0$.
Then $n=1$ so
$\eta =\eta_1 \in \lbrace p, p^{-1}\rbrace$.
The elements $p$, $p^{-1}$ are not in~$Z(\Lambda)$, for a contradiction.
Therefore $\ell \neq 0$.
From the nature of the matrices~$p$,~$s$ in Def\/inition~\ref{def:ps},
we may
view $\eta $ as a polynomial in $\lambda$
that has
coef\/f\/icients in
${\rm {Mat}}_2(\mathbb F)$ and degree at most $\ell$.
Call this polynomial~$f$.
We claim that the degree of~$f$ is exactly~$\ell$.
To prove the claim,
write
\begin{gather*}
s = s_0+ s_1 \lambda,   \qquad
s_0= \left(
\begin{array}{  c c}
0 & 1  \\
0 & 0
\end{array}
\right),
\qquad
s_1= \left(
\begin{array}{  c c}
0 & 0  \\
1 & 0
\end{array}
\right).
\end{gather*}
Let $m
\in {\rm {Mat}}_2(\mathbb F)$ denote the coef\/f\/icient of
$\lambda^\ell$ in $f$.
The matrix $m$ is obtained from
$\eta_1 \eta_2\cdots \eta_n$ by
replacing each occurrence of $s$ by $s_1$.
Using $s_1p s_1=-s_1$ and
 $s_1p^{-1} s_1=s_1$  we f\/ind
$m\in \lbrace \pm p^i s_1 p^j |
-1 \leq i,j\leq 1\rbrace $. The matrix $p$ is invertible
and $s_1\not=0$ so
$p^i s_1 p^j\not=0$ for $-1 \leq i,j \leq 1$.
Therefore $m\not=0$ and the claim is proved.
Let $\kappa $ denote the $(1,1)$-entry of the matrix $\eta $.
Then $\eta=\kappa I$ since $\eta \in Z(\Lambda )$.
In the equation
$\eta_1 \eta_2\cdots \eta_n=\kappa I$ take the determinant of
each side and use
Lemma~\ref{lem:4}$(i)$
to get $(-\lambda)^\ell = \kappa^2$.
Therefore $\ell $ is even and
$\kappa=\pm \lambda^{\ell/2}$.
Now $\eta =\pm \lambda^{\ell/2} I$, so
the above polynomial $f$ has degree $\ell/2$.
But $\ell \not=0$ so $\ell > \ell/2$
for a contradiction.
the result follows.
\end{proof}

We now display an algebra homomorphism $\Delta \to \Lambda $.

\begin{Lemma}
\label{prop:uhomoh}
There exists a unique $\mathbb F$-algebra homomorphism
$\pi:\Delta \to \Lambda $
 that sends
\begin{gather*}
A \mapsto q \mathcal A+q^{-1}{\mathcal A}^{-1},
 \qquad
B \mapsto q \mathcal B+q^{-1}{\mathcal B}^{-1},
\qquad
C \mapsto q\mathcal C+q^{-1}{\mathcal C}^{-1}.
\end{gather*}
The homomorphism $\pi$ sends
\begin{gather}
 \alpha \mapsto \nu I,
 \qquad
\beta \mapsto  \nu I,
 \qquad
\gamma \mapsto \nu I,
\label{eq:pisend}
\end{gather}
where $\nu= (q^2+q^{-2})\mu + \mu^2$.
\end{Lemma}

\begin{proof}
Def\/ine
\begin{gather*}
A^\vee = q\mathcal A+q^{-1}{\mathcal A}^{-1},
\qquad
B^\vee =q\mathcal B+q^{-1}{\mathcal B}^{-1},
\qquad
C^\vee = q\mathcal C+q^{-1}{\mathcal C}^{-1}.
\end{gather*}
By Lemma~\ref{lem:xyz}
and Lemma~\ref{lem:multtable2},
\begin{gather}
(q+q^{-1})A^\vee+ \frac{qB^\vee C^\vee-q^{-1}C^\vee B^\vee}{q-q^{-1}}
  = \nu I,
\label{eq:vee1}
\\
(q+q^{-1})B^\vee+
\frac{qC^\vee A^\vee -q^{-1}A^\vee C^\vee }{q-q^{-1}}
 =  \nu I,
\label{eq:vee2}
\\
(q+q^{-1})C^\vee +
\frac{qA^\vee B^\vee -q^{-1}B^\vee A^\vee }{q-q^{-1}}
 =  \nu I.
\label{eq:vee3}
\end{gather}
By
(\ref{eq:vee1})--(\ref{eq:vee3}) and
since
 $\nu I$ is central, the elements
$A^\vee$, $B^\vee$, $C^\vee$ satisfy the
def\/ining relations for~$\Delta $ from
Def\/inition~\ref{def:uaw}.
Therefore
the homomorphism $\pi$ exists. The homomorphism $\pi$
is unique since $A$, $B$, $C$ generate~$\Delta $. Line
(\ref{eq:pisend})
follows from
Def\/inition~\ref{def:abc} and
(\ref{eq:vee1})--(\ref{eq:vee3}).
\end{proof}



\begin{Lemma}
\label{lem:comdiagf}
For $g \in
{\rm  {PSL}}_2(\mathbb Z)$
 the following diagram commutes:
\[
\begin{CD}
\Delta  @>\pi>>
         \Lambda  \\
	  @Vg VV                     @VVg V \\
	 \Delta
	  @>\pi>>
	          \Lambda
		    \end{CD}
		    \]

\end{Lemma}

\begin{proof}
The elements
$\rho$, $\sigma$ form a generating set for
${\rm  {PSL}}_2(\mathbb Z)$; without loss
we may assume that~$g$ is contained in this set.
By Theorem~\ref{thm:try}
the action of~$\rho$ on~$\Delta $
cyclically permutes~$A$,~$B$,~$C$.
By Lemma~\ref{lem:4}$(iii)$
the action of $\rho$ on $\Lambda $
cyclically permutes
$\mathcal A$, $\mathcal B$, $\mathcal C$.
By Theorem~\ref{thm:try}
the action of~$\sigma$ on~$\Delta $ swaps $A$, $B$ and f\/ixes
$\gamma$. By Lemma~\ref{lem:4}$(iv)$ and the construction,
the action of $\sigma $ on~$\Lambda $ swaps $\mathcal A$, $\mathcal B$
and f\/ixes $ I$.
The diagram commutes by these comments
and Lemma~\ref{prop:uhomoh}.
\end{proof}

\begin{Theorem}
\label{thm:faith}
The action of
${\rm  {PSL}}_2(\mathbb Z)$ on $\Delta $ is faithful.
\end{Theorem}

\begin{proof}
Let $g$ denote an element of
${\rm  {PSL}}_2(\mathbb Z)$ that f\/ixes everything in
$\Delta $.
We show that $g=1$.
By Lemma~\ref{lem:conj} and since $\rho$, $\sigma$ generate
${\rm  {PSL}}_2(\mathbb Z)$, there exists an invertible
$\xi \in \Lambda $ such that $g(\eta)=\xi \eta \xi^{-1}$ for
all $\eta  \in \Lambda $. By assumption
$g$ f\/ixes the element $A$ of $\Delta $. Under
the homomomorphim
$\pi : \Delta  \to \Lambda $ the image of $A$ is
$q\mathcal A+q^{-1}{\mathcal A}^{-1}$,
so $g$ f\/ixes
$q\mathcal A+q^{-1}{\mathcal A}^{-1}$ in view of Lemma~\ref{lem:comdiagf}. Therefore $\xi $ commutes
with
$q\mathcal A+q^{-1}{\mathcal A}^{-1}$.
Recall $\mathcal A+{\mathcal A}^{-1} =\mu I$ by
Lemma~\ref{lem:xyz}, so $\xi$ commutes with
$\mathcal A+{\mathcal A}^{-1}$.
By these comments and $q^2\not=1$ we f\/ind
$\xi$ commutes with $\mathcal A$. By a similar
argument $\xi$ commutes with $\mathcal B$. Now
$\xi \in Z(\Lambda )$ by
Lemma~\ref{lem:whatittakes}.
Consequently $g$ f\/ixes everything in $\Lambda $, so $g=1$
 by Lemma~\ref{lem:faith1}.
\end{proof}

\begin{thebibliography}{99}

\footnotesize\itemsep=0pt
\bibitem{alpern}
Alperin R.C.,
Notes: ${\rm  {PSL}}_2(\mathbb Z)= \mathbb Z_2 \star \mathbb Z_3$,
\href{http://dx.doi.org/10.2307/2324963}{{\it Amer. Math. Monthly}} {\bf 100} (1993), 385--386.
\end{thebibliography}
\end{document}
